\section{Metodología y enofoque propuesto}

\subsection{Enfoque técnico}
Descripción de la propuesta de solución, indicando las técnicas que se utilizarán, con su justificación (cómo se alinea con los objetivos). Indicar la metodología de desarrollo de la solución y un diagrama de la solución propuesta.

This project will begin with a thorough analysis of the dataset, presented in the following section. During this stage, we will identity it's structure, composition and possible pre-processing filters that can be applied to enhance the training data.

Then, we will begin with the design of an ML/DL architecture based on the foundational models established in EgoVLP, LAVILA and JPoSE. Starting off from foundational models proves useful in fine-tuning and will enable an exploration of the impact of processing in multimodal retrieval. 

To implement and experiment this architecture, we will use modern and robust MLOps pipeline that pulls weights from AVION ViT-L/14, SMS-Loss (2024) and CR-CLIP (2025). Experiment management will be done with \href{https://hydra.cc/}{HYDRA} and composable YAML files.

This pipeline will enable experimentation and clear analysis of metrics (nDCG and mAP) in order to select and iterate on the best results.

\subsection{Herramientas}

The following libraries will be used:

\begin{lstlisting}
    PyTorch, torchvision, torchaudio (CUDA 12.1)
    transformers, timm, einops
    hydra-core, omegaconf, pydantic
    wandb, tensorboard
    matplotlib, seaborn, plotly
\end{lstlisting}
